\section[Application to the construction of formal schemes...]{Application to the construction of formal schemes
and smooth ordinary schemes over a complete local ring~$A$}
\label{III.7}
%
%
The results of the preceding~\No sometimes make it possible to
prove the existence
%
% original p. 83
%
of an $\goth{m}$-adic formal scheme over such a ring,
reducing to a given smooth scheme~$X_0$ over~$k$.
Let us distinguish two cases:

a) \emph{$A$ is ``of equal characteristic''} (this is the case
in particular if~$k$ has characteristic~$0$). Then one knows
that there exists a \emph{subfield of representatives of}~$A$,
\ie a subfield~$k'$ such that $A\to k$ induces an isomorphism~$k'\isomto
k$. \emph{Then there even exists a smooth ordinary scheme
over~$A$ reducing to~$X_0$}, namely~$X=X_0\otimes_kA$, with~$A$
regarded as an algebra over~$k$ by means of the homomorphism
$k\to k'\to A$ defined by~$k'$. It should nevertheless be noted
that this construction is not ``natural''; it is easy to convince
oneself (already in the case where~$A=k[t]/(t^2)$, the algebra of
dual numbers) that another lifting homomorphism~$k\to A$ (defined
in this instance by an absolute derivation of~$k$ into itself)
defines an~$X'$ over~$A$ which in general \emph{is not isomorphic
to~$X$} (if~$\H^1(X_0,\goth{g}_{X_0/k})\neq0$). It would moreover
be interesting to study (for~$k$ of characteristic~$0$, or
non-perfect and of characteristic~$p>0$) which smooth~$X$ over~$A$
are obtained in this way, and under what condition two homomorphisms
$k\to A$ define isomorphic $A$-schemes. Nevertheless, the existence
of~$k'$ suffices to entail that the first obstruction to lifting
$X_0$, which is
in~$\H^2(X_0,\goth{g}_{X_0/k})\otimes_k\goth{m}/\goth{m}^2$, is
necessarily zero. Of course, when one has then lifted~$X_0$ to~$X_1$
smooth over~$A/\goth{m}^2$, the new obstruction to constructing~$X_2$
will in general not be zero: it will be a function of a variable
element in a certain principal homogeneous space
under~$\H^1(X_0,\cal{G}_0)\otimes\goth{m}/\goth{m}^2$ and lies
in~$\H^2(X_0,\cal{G}_0)\otimes\goth{m}^2/\goth{m}^3$: it would
be appropriate to study the situation in detail\footnote{It is
no doubt described by the Kodaira--Spencer bracket operation (\cf
Cartan Seminar, 1960/61, Exp\ptbl 4).}.

b) \emph{$A$ is of unequal characteristics.} In this case, we
know nothing (except if by chance~$\H^2(X_0,\goth{g}_{X_0/k})=0$,
in which case one can construct a simple $\goth{m}$-adic formal
scheme over~$A$ reducing to~$k$). Even if~$A=\ZZ/p^2\ZZ$ and~$X_0$
is an ``abelian'' scheme of dimension~$2$, one does not know whether
it can be lifted to an~$X=X_1$ smooth over~$A$\footnote{This has
now been proved, \cf note~\ref{III.6.6.p24} page~\pageref{III.6.6.p24}.},
on the other hand, one has no example of an~$X_0$ which one has
proved does not come from an ordinary scheme~$X$ smooth over~$A$.
(I have the impression that such an example must exist, with~$X_0$
a projective surface)\footnote{\label{nIII.3}Such an example has since been
constructed by
J-P\ptbl Serre (Proc.\ Nat.\ Acad.\ Sc.\ USA \textbf{47} (1961),
\No 1, p\ptbl108--109), at least for certain
dimensions. D\ptbl Mumford gave an example (unpublished) with
an algebraic \emph{surface}.}. Let us simply point out that,
according to Cohen's theorem, there exists a Cohen $p$-ring~$B$
with residue field~$k$ and a homomorphism~$B\to A$ inducing the
identity isomorphism on the residue fields; consequently, the
``strongest'' lifting result would be obtained by taking~$A$
%
% original p. 84
%
a Cohen $p$-ring: if there exists a solution (ordinary or formal)
over such a ring, there exists one over every complete local ring
with residue field~$k$. In particular, since for a Cohen $p$-ring
$\goth{m}/\goth{m}^2$ is canonically identified with~$k$, one sees
that \emph{for every smooth scheme~$X_0$ over a field~$k$ of
characteristic~$p>0$, there exists a cohomology class
in~$\H^2(X_0,\goth{g}_{X_0/k})$} which is the first obstruction to
lifting~$X_0$ to a smooth scheme over a Cohen $p$-ring; one does
not know whether it can be nonzero\footnote{It can be nonzero,
as indicated in note~\ref{nIII.3} on page~\pageref{nIII.3}.}.

Even if one succeeds step by step in constructing the~$X_n$
reducing to~$X_0$, this generally gives only a \emph{formal}
scheme~$\goth{X}$ smooth over~$A$, reducing to~$X_0$. When~$X_0$
is proper over~$A$, the question remains whether~$\goth{X}$ is in
fact algebraizable, in order to obtain an \emph{ordinary} scheme
proper over~$A$ and simple over~$A$, reducing to~$X_0$. The only
known criterion (mentioned in the Bourbaki Seminar, and appearing
in the Elements, Chapter~III, 4.7.1) is the following: if~$\goth{X}$
is proper over~$A$, and if~$\cal{L}$ is an invertible sheaf on~$\goth{X}$
such that the sheaf induced~$\cal{L}_0$ on~$X_0$ is ample (\ie a suitable
tensor power~$\cal{L}_0^{\otimes n}$, $n>0$, comes from a projective
embedding of~$X_0$), then there exists a scheme~$X$ projective over~$A$,
and an ample invertible sheaf on~$X$, such that~$(\goth{X},\cal{L})$
is obtained from it by $\goth{m}$-adic completion. This therefore
leads us, given a locally free sheaf~$\cal{E}_0$ on~$X_0$ (which we
shall choose ample and invertible for our purpose), to extend it to
a locally free sheaf~$\cal{E}$ on~$\goth{X}$. For this, we are
reduced to constructing step by step locally free sheaves~$\cal{E}_n$
on the~$X_n$. The discussion is entirely analogous to that in
\No\ref{III.6}, (\cf remark~\ref{III.6.4}), the essential role being
played by the \emph{sheaf of automorphisms} of an~$\cal{E}_{n+1}$
which induce the identity on~$\cal{E}_n$: one immediately shows that
this sheaf is identified with
\begin{equation*}
\label{eq:III.7.*}\tag{$*$}{}
\cal{G}= \SheafHom_{\cal{O}_{X_0}}(\cal{E}_0,
\cal{E}_0\otimes\gr_{\cal{I}}^{n+1}(\cal{O}_X))=
\SheafHom_{\cal{O}_{X_0}}(\cal{E}_0,
\cal{E}_0)\otimes\gr_{\cal{I}}^{n+1}(\cal{O}_X)
\end{equation*}
which is again a sheaf of commutative groups. We find:

\begin{proposition}\label{III.7.1}
Let~$S$ be a prescheme endowed with a quasi-coherent sheaf of
ideals~$\cal{I}$, $X$ a prescheme over~$S$, $S_n$ the subprescheme
of~$S$ defined by~$\cal{I}^{n+1}$, and~$X_n=X\times_SS_n$ (for every
integer~$n$). Let~$\cal{E}_n$ be a locally free sheaf
%
% original p. 85
%
on~$X_n$; we propose to extend it to a locally free sheaf~$\cal{E}_{n+1}$
on~$X_{n+1}$. Then~$\cal{E}_n$ defines a canonical obstruction class
in~$\H^2(X_0,\cal{G})$, where~$\cal{G}$ is the quasi-coherent sheaf
given by the formula above, whose vanishing is necessary and
sufficient for the existence of an~$\cal{E}_{n+1}$ extending~$\cal{E}_n$.
If this class is zero, then the set of classes, up to isomorphism
(inducing the identity on~$\cal{E}_n$), of solutions~$\cal{E}_{n+1}$,
is a principal homogeneous space under~$\H^1(X_0,\cal{G})$.
\end{proposition}

This proposition gives rise to the usual corollaries. Let us point
out only that if~$X$ is \emph{flat} over~$S$, then one can write
$$
\cal{G}= \SheafHom_{\cal{O}_{X_0}}(\cal{E}_0, \cal{E}_0)
\otimes_{\cal{O}_{S_0}} \gr_{\cal{I}}^{n+1}(\cal{O}_S)
$$
whence, if~$S$ is affine of ring~$A$ and if the~$\cal{I}^n/\cal{I}^{n+1}$
are locally free, the sufficient condition
$$
\H^2(X_0,\cal{G}_0)=0 \qquad\text{with}\qquad
\cal{G}_0=\SheafHom_{\cal{O}_{X_0}}(\cal{E}_0, \cal{E}_0)
$$
for the existence of an~$\cal{E}_{n+1}$, and hence step by step for
the existence of successive extensions~$\cal{E}_m$ ($m=n,n+1,$
etc.).

Returning to the initial situation, we therefore find:

\begin{proposition}\label{III.7.2}
Let~$A$ be a complete local ring,~$\goth{X}$ a proper and flat formal
scheme over~$A$, such that~$X_0$ is projective and
$\H^2(X_0,\cal{O}_{X_0})=0$. Then there exists a scheme~$X$
projective over~$A$ whose $\goth{m}$-adic formal completion is
isomorphic to~$\goth{X}$.
\end{proposition}

Combining this with~\ref{III.6.10}, we find:

\begin{theorem}\label{III.7.3}
Let~$A$ be a complete local ring with residue field~$k$, $X_0$ a
projective and smooth scheme over~$k$, such that
$$
\H^2(X_0,\goth{g}_{X_0/k})=\H^2(X_0,\cal{O}_{X_0})=0
$$
Then there exists a smooth and projective scheme~$X$ over~$A$,
reducing to~$X_0$.
\end{theorem}

More generally, if one is given an~$X_n$ smooth over~$A_n=A/\goth{m}^{n+1}$
reducing to~$X_0$, then there exists an~$X$ smooth and proper over~$A$
and an isomorphism~$X\otimes_AA_n=X_n$.

\begin{corollary}\label{III.7.4}
Every smooth and proper curve over~$k$ is obtained by reduction from
a smooth and proper curve over~$A$.
\end{corollary}

It is this result which will be the essential tool (together with the
existence theorem for sheaves in Formal Geometry) for studying the
fundamental group of~$X_0$ by transcendental means.
%
% original p. 86
%
